% Kandydat restrukturyzacji; nie jest aktywnym rozdziałem Atlasu.
\section{Interpolacja rzeczywista i twierdzenie Marcinkiewicza}
\providecommand{\atlasref}[2]{#2}

\begin{lemat}[Interpolacja Marcinkiewicza pomiędzy słabym $L^1$ i silnym $L^2$]
\label{lem:dt-cz-interp}
Niech $T$ będzie ograniczonym operatorem liniowym
na $L^2(\R^d)$, spełniającym oszacowanie słabego typu
$(1,1)$ na $L^1\cap L^2$ ze stałą $C_1$.
Dla każdego $1<p<2$ jego ograniczenie do
$L^1\cap L^2\cap L^p$ spełnia
$\|Tf\|_p\le C_p\|f\|_p$, gdzie $C_p$ zależy również
od stałych obu oszacowań. Ma ono jednoznaczne
ograniczone przedłużenie na $L^p$.
\end{lemat}
\begin{proof}
Dla $f\in L^1\cap L^2\cap L^p$ i $t>0$ kładziemy
$f_t^{>}=f\chi_{\{|f|>t\}}$ i
$f_t^{\le}=f\chi_{\{|f|\le t\}}$.
Obie funkcje należą do $L^1\cap L^2$.
Z liniowości i nierówności trójkąta wynika
$\{|Tf|>2t\}\subset\{|Tf_t^{>}|>t\}
\cup\{|Tf_t^{\le}|>t\}$ z dokładnością do zbioru
miary zero. Oznaczając $\|T\|_{L^2\to L^2}=C_2$,
otrzymujemy
\[
 \lambda_d(\{|Tf|>2t\})
 \le \frac{C_1}{t}\int_{\{|f|>t\}}|f|\,\mathrm d\lambda_d
 +\frac{C_2^2}{t^2}\int_{\{|f|\le t\}}|f|^2\,\mathrm d\lambda_d .
\]
Tożsamość
$|Tf(x)|^p=\int_0^{|Tf(x)|}p s^{p-1}\,\mathrm d\lambda_1(s)$
i twierdzenie Tonellego dają tożsamość warstwową
\[
 \|Tf\|_p^p=p\int_0^\infty s^{p-1}
 \lambda_d(\{|Tf|>s\})\,\mathrm d\lambda_1(s)
\]
również wtedy, gdy początkowo obie strony są nieskończone.
Po podstawieniu $s=2t$ i użyciu Tonellego otrzymujemy
\[
 \begin{aligned}
 \|Tf\|_p^p
 &\le p2^p C_1\int_{\{|f|>0\}} |f(x)|
       \int_0^{|f(x)|}t^{p-2}\,\mathrm d\lambda_1(t)
       \,\mathrm d\lambda_d(x)\\
 &\quad+p2^p C_2^2\int_{\{|f|>0\}} |f(x)|^2
       \int_{|f(x)|}^\infty t^{p-3}\,\mathrm d\lambda_1(t)
       \,\mathrm d\lambda_d(x)\\
 &=p2^p\left(\frac{C_1}{p-1}+\frac{C_2^2}{2-p}\right)
   \|f\|_p^p .
 \end{aligned}
\]
Punkty $f=0$ pomijamy, gdyż pierwotne całki po $t$
są w nich zerowe. Warunki $p>1$ i $p<2$ odpowiadają
zbieżności obu całek wewnętrznych.
Gęstość $C_c^\infty$ w $L^p$
(twierdzenie~\ref{thm:dt-test-density}) i zupełność $L^p$
dają jednoznaczne przedłużenie.
Jest ono zgodne z wyjściowym operatorem $T$ na
$L^2\cap L^p$: dla $f$ z tej przestrzeni funkcje
$f_n=f\chi_{B_n(0)}\chi_{\{|f|\le n\}}$
należą do $L^1\cap L^2\cap L^p$ i zbiegają do $f$
w obu normach przez zbieżność zmajoryzowaną.
Ich obrazy zbiegają w $L^2$ do $Tf$ oraz w $L^p$
do wartości przedłużenia. Z obu zbieżności wybieramy
wspólny podciąg zbieżny prawie wszędzie, więc granice
są równe prawie wszędzie.
\end{proof}
